% Section 4: The Constitutional Enclosure Theorem
% IEEE Computer — "Silent Decisions Are Type Errors"
% Authors: Soares et al., 2026

\section{The Constitutional Enclosure Theorem}
\label{sec:theorem}

This section formalises the paper's bounded claim: an external Safe Rust
caller cannot obtain a BTV \texttt{Verdict} through the crate's public API
without supplying the required evidence and compliance tokens.  This is an
API-enclosure result, not a theorem that every consequential action in a
surrounding system is routed through BTV.

% ----------------------------------------------------------------
\subsection{Definitions}
\label{sec:definitions}

\begin{definition}[Materialized Verdict]
\label{def:materialized}
A verdict $V$ is \emph{materialized} if and only if there exist an
evidence token $E$, a compliance token $C$, a decision outcome
$d \in \{\mathit{Allow},\, \mathit{Deny}\}$, and an explanation string
$x$ such that
\[
  V = \texttt{Verdict::new}(E,\, C,\, d,\, x).
\]
The function \texttt{Verdict::new} is the \emph{sole} constructor of
the type \texttt{Verdict}; no other term of that type is well-typed in
Safe Rust.
\end{definition}

\begin{definition}[Silent Decision]
\label{def:silent}
For a program $P$, let $\mathcal{S}_{P}$ denote the set of
\emph{unevidenced BTV verdicts}: terms of type \texttt{Verdict} exposed
through the crate's public API without construction from one
\texttt{EvidenceToken} and one authority-issued \texttt{ComplianceToken}.
This definition is deliberately confined to the BTV type perimeter; an
external action that bypasses \texttt{Verdict} is outside the set.
\end{definition}

% ----------------------------------------------------------------
\subsection{Axioms}
\label{sec:axioms}

\begin{axiom}[Resource Introduction, $\otimes$-I]
\label{ax:intro}
If $\Gamma \vdash E : \texttt{EvidenceToken}$ and
$\Delta \vdash C : \texttt{ComplianceToken}$, then
\[
  \Gamma,\, \Delta \;\vdash\; \texttt{Verdict::new}(E,\, C,\, d,\, x)
  : \texttt{Verdict}.
\]
The contexts $\Gamma$ and $\Delta$ are \emph{disjoint}: linear logic
forbids sharing a resource between both premises.  In Rust this is
enforced by the move semantics of \texttt{Verdict::new}, which takes
$E$ and $C$ by value.  Once passed to the constructor, neither token
is accessible to the caller.
\end{axiom}

\begin{axiom}[Resource Uniqueness and Drop Diagnostics]
\label{ax:unique}
\texttt{EvidenceToken} is annotated \texttt{\#[must\_use]} and does
not implement \texttt{Clone} or \texttt{Copy}.  Consequently:
\begin{enumerate}
  \item \emph{Contraction} is prohibited: a token cannot be
    duplicated, preventing the same evidence from being
    associated with more than one verdict.
  \item Ignoring a token-producing expression emits an
    \texttt{unused\_must\_use} diagnostic, treated as an error in this
    crate.  Explicit discard remains possible because Rust is affine;
    this lint is a guardrail, not an anti-weakening proof.
\end{enumerate}
The method \texttt{consume()} is the unique destructor; it transfers
the inner \texttt{Blake3Hash} to the \texttt{Verdict} constructor and
is declared \texttt{pub(crate)}, making it inaccessible outside the
defining crate.
\end{axiom}

\begin{axiom}[Encapsulation Boundary]
\label{ax:encap}
Let $\Gamma_{\mathit{ext}}$ denote any typing context external to the
defining crate (i.e., any context in a crate that depends on
\texttt{silent-decisions-proof} but is not the crate itself).
The following are ill-typed in $\Gamma_{\mathit{ext}}$:
\begin{enumerate}
  \item Struct-literal construction of \texttt{Verdict}---all fields
    are private (compiler error \texttt{E0451}).
  \item Direct invocation of \texttt{EvidenceToken::consume()}:
    the method is \texttt{pub(crate)} (compiler error \texttt{E0624}).
  \item Construction of \texttt{Blake3Hash} via its tuple-struct
    field---the field is private (compiler error \texttt{E0423}).
\end{enumerate}
Together, these restrictions close the encapsulation perimeter: the
only well-typed path to a \texttt{Verdict} in $\Gamma_{\mathit{ext}}$
is through \texttt{Verdict::new}, which enforces
Axiom~\ref{ax:intro}.
\end{axiom}

% ----------------------------------------------------------------
\subsection{Threat Model}
\label{sec:threat}

The theorem holds under the \emph{Safe Rust} memory model.
The \texttt{unsafe} keyword grants the programmer the ability to
bypass borrow-checker and visibility rules; consequently, code that
uses \texttt{unsafe} is explicitly out of scope.

First-party crates forbid \texttt{unsafe}; production deployments should
also inventory dependency usage with \texttt{cargo geiger}~\cite{cargogeiger}.
That tool reports unsafe constructs but does not prove their soundness or
guarantee that the theorem composes across the entire dependency graph.

Macro-generated code is in scope: procedural macros run before the
type checker and cannot produce terms that circumvent visibility rules
without themselves using \texttt{unsafe}.

% ----------------------------------------------------------------
\subsection{Proof of the Constitutional Enclosure Theorem}
\label{sec:proof}

\begin{theorem}[Constitutional Enclosure]
\label{thm:enclosure}
Assume the BTV crate implementation and its public encapsulation are
preserved.  No external Safe Rust caller can obtain a term of type
\texttt{Verdict} through that public API without the constructor consuming
one \texttt{EvidenceToken} and one authority-issued
\texttt{ComplianceToken}.  Equivalently, $\mathcal{S}_{P}=\varnothing$
within this API perimeter.
\end{theorem}

\begin{proof}
We analyze the public construction and extraction paths closed by the
crate's encapsulation assumptions.

\medskip
\noindent\textbf{Case~A: Struct-literal construction.}
An adversary attempts to write
\begin{verbatim}
  Verdict { decision: Decision::Allow, .. }
\end{verbatim}
from a crate external to \texttt{silent-decisions-proof}.
By Axiom~\ref{ax:encap}(1), all fields of \texttt{Verdict} are
private.  The compiler rejects this expression with
\texttt{error[E0451]: field `evidence\_id` of struct `Verdict` is
private}.  This is verified by the compile-fail test
\texttt{tests/ui/verdict\_struct\_literal.rs}.

\medskip
\noindent\textbf{Case~B: External invocation of \texttt{consume()}.}
An adversary attempts to call
\texttt{EvidenceToken::consume()} directly from an external crate,
bypassing \texttt{Verdict::new} to extract a \texttt{Blake3Hash} and
assemble a \texttt{Verdict} fragment.
By Axiom~\ref{ax:encap}(2), \texttt{consume()} is declared
\texttt{pub(crate)}.  The compiler rejects the call with
\texttt{error[E0624]: method `consume` is private}.
This is verified by the compile-fail test
\texttt{tests/ui/external\_consume\_call.rs}.

\medskip
\noindent\textbf{Guardrail~C: Discard of \texttt{EvidenceToken}.}
A developer issues a decision without providing evidence by
constructing an \texttt{EvidenceToken} and immediately dropping it
(or ignoring the return value of a function that returns one).
By Axiom~\ref{ax:unique}, ignoring a token-producing expression triggers
a diagnostic promoted to an error in this crate.  Explicit
\texttt{let \_ = expr}, \texttt{drop}, or \texttt{mem::forget} remains
possible.  None produces a \texttt{Verdict}; the formal closure instead
comes from Axiom~\ref{ax:encap} (private fields and restricted visibility).
This is verified by the compile-fail test
\texttt{tests/ui/dropped\_evidence\_token.rs}.

\medskip
By the Encapsulation Boundary axiom, no term of type \texttt{Verdict}
can be obtained in $\Gamma_{\mathit{ext}}$ without passing through
\texttt{Verdict::new}.  By Definition~\ref{def:materialized}, every
such term is a Materialized Verdict; therefore
$\mathcal{S}_{P}=\varnothing$ within the stated perimeter.
\end{proof}

% ----------------------------------------------------------------
\subsection{Corollary: Fail-Secure by Construction}
\label{sec:corollary}

\begin{corollary}[Fail-Secure]
\label{cor:failsecure}
In any Safe Rust program satisfying the conditions of
Theorem~\ref{thm:enclosure}, the code path that handles a runtime
failure to produce evidence (e.g., a database timeout, a missing
context field) cannot silently issue a \texttt{Verdict}---the type
system ensures at compile time that every branch leading to a
\texttt{Verdict} must supply a valid \texttt{EvidenceToken}.
The protected path \emph{fails securely}: it cannot return a BTV verdict
that lacks an evidence record.
\end{corollary}

\begin{proof}
If the evidence-production path is absent or returns an error type
(e.g., \texttt{Result<EvidenceToken, E>}), the only well-typed
continuation that produces a \texttt{Verdict} requires unwrapping
the result and handling the error branch explicitly.  The error
branch cannot produce a \texttt{Verdict} without evidence; it must
either propagate the error upward or halt.  By
Theorem~\ref{thm:enclosure}, any attempt to construct a
\texttt{Verdict} in the error branch without a valid
\texttt{EvidenceToken} is ill-typed.
\end{proof}

The Fail-Secure Corollary supports review and explanation under
LGPD Art.~20 and EU AI Act Arts.~14 and~86 by making the required
artifacts structural on the protected BTV path rather than an optional
post-hoc assertion.

% ----------------------------------------------------------------
\subsection{Corollary 4.8: Human Escalation as the Only Well-Typed Alternative}
\label{sec:escalated}

\begin{corollary}[Human Escalation]
\label{cor:escalated}
Under the conditions of Theorem~\ref{thm:enclosure}, when an
\texttt{EvidenceToken} cannot be produced (e.g., model timeout,
unavailable context), the API provides a well-typed alternative to error
propagation or halt: constructing an \texttt{EscalatedVerdict}
via a linearly-consumed \texttt{OperatorToken}.
Silent fallback to an unevidenced automated verdict is ill-typed.
Formally, $O \multimap V_{\mathrm{esc}}$.

This structure aligns with the human oversight requirement of
EU AI Act Art.~14: the system cannot produce a consequential
decision-equivalent record without binding exactly one named
human operator, whose identity and authorisation are protected by an
HMAC integrity seal under the configured key-management model.
\end{corollary}

\begin{proof}
\texttt{EscalatedVerdict} is a sibling type to \texttt{Verdict}
and does not affect the proof of Theorem~\ref{thm:enclosure}.
\texttt{std::mem::size\_of::\textless{}EscalatedVerdict\textgreater{}()}
measures $160$ bytes versus $184$ for \texttt{Verdict}
(x86-64, rustc~1.98.1; Section~\ref{sec:results}) --- the automated
verdict is the larger type because it carries the signed compliance
token's 32-byte signature.
Its sole constructor \texttt{EscalatedVerdict::new} requires an
\texttt{OperatorToken} by value, following the same
$\otimes$-I pattern as Axiom~\ref{ax:intro}.
\texttt{OperatorToken} carries the same linear-resource invariants
as \texttt{EvidenceToken}: it is \texttt{\#[must\_use]},
not \texttt{Clone}/\texttt{Copy}, and its \texttt{consume()}
destructor is \texttt{pub(crate)}.
The four compile-fail tests
\texttt{escalated\_struct\_literal.rs},
\texttt{escalated\_token\_reuse.rs},
\texttt{escalated\_operator\_token\_drop.rs}, and
\texttt{escalated\_consume\_external.rs}
exercise the corresponding encapsulation and ownership cases
of Theorem~\ref{thm:enclosure}) that no \texttt{EscalatedVerdict}
can be forged, duplicated, or silently dropped.
\end{proof}

% ----------------------------------------------------------------
\subsection{Polyglot Deployment: The FFI Trust Boundary}
\label{sec:ffiboundary}

Most production ML pipelines are polyglot: orchestration logic (agent
frameworks, model servers, API gateways) runs in Python and other
languages, while BTV's linear-type discipline is enforced by the Rust
compiler.
The PyO3 gateway (\texttt{btv-python}) is the sole authorized boundary
between these worlds: a Python caller can invoke
\texttt{issue\_verdict}, but cannot inspect, clone, or reconstruct a
\texttt{Verdict}---it receives only an opaque, integrity-sealed
\texttt{SealedVerdict} handle exposing read-only accessors and a
\texttt{verify\_integrity} check.
Theorem~\ref{thm:enclosure} is therefore enforced strictly within the
\texttt{btv-core} and \texttt{btv-python} crates; it does not, and
cannot, extend to the untrusted orchestrator itself.
Consequently, the enforcement burden shifts downstream: the effector
that materialises a decision's real-world consequence (the
credit-denial API, the database write that blocks a transaction) must
independently require and verify the sealed verdict handle before
acting, and must reject any request that bypasses the gateway.
The Rust enclave protects itself; it does not, by itself, protect the
system that surrounds it.
